\documentclass[10pt]{article}
\usepackage[T1]{fontenc}
\usepackage[utf8]{inputenc}
\usepackage[french]{babel}
\usepackage{amsmath,amssymb,mathrsfs,mathpazo}
\usepackage{geometry}
\geometry{a4paper,margin=2cm}
\usepackage{graphicx,array,tabularx,tabularray,colortbl,multirow,multicol,booktabs,enumitem,pifont,ulem,fancybox,diagbox,cancel}
\usepackage{pstricks,pstricks-add,pst-plot,pst-tree,pst-node,pst-text,pst-func,pst-math,pst-bspline,pst-3dplot,pst-eucl,pst-solides3d}
\let\pscalculate\undefined
\usepackage{pst-bezier}
\usepackage{tikz,pgfplots,tkz-tab}
\usetikzlibrary{arrows,arrows.meta,patterns,calc,positioning,trees,decorations.markings,shapes.geometric,shapes.misc,shapes.arrows,intersections,decorations.pathmorphing,decorations.pathreplacing,backgrounds,shadows,babel}
\pgfplotsset{compat=1.18}
\usepgfplotslibrary{fillbetween}
\usepackage[np]{numprint}
\usepackage{listings,makecell,pas-tableur,twoopt}
\let\barre\undefined
\usepackage[pstricks]{tablvar}
\definecolor{darkgreen}{RGB}{0,100,0}
\providecommand{\red}{\color{red}}
\providecommand{\blue}{\color{blue}}
\providecommand{\green}{\color{green}}
\providecommand{\black}{\color{black}}
\providecommand{\white}{\color{white}}
\providecommand{\gray}{\color{gray}}
\providecommand{\cyan}{\color{cyan}}
\providecommand{\magenta}{\color{magenta}}
\providecommand{\yellow}{\color{yellow}}
\providecommand{\violet}{\color{violet}}
\providecommand{\orange}{\color{orange}}
\providecommand{\euro}{\texteuro}
\providecommand{\EUR}{\texteuro}
\providecommand{\ieme}{\textsuperscript{e}}
\providecommand{\degres}{\ensuremath{^\circ}}
\providecommand{\degree}{\ensuremath{^\circ}}
\providecommand{\vv}[1]{\overrightarrow{#1}}
\providecommand{\Coord}[1]{#1}
\providecommand{\decosix}{}
\providecommand{\starredbullet}{\ensuremath{\bullet}}
\pagestyle{empty}
\setlength{\parindent}{0pt}
\usepackage{adjustbox}\DeclareUnicodeCharacter{2236}{:}\DeclareUnicodeCharacter{2190}{\ensuremath{\leftarrow}}
\providecommand{\R}{\mathbb{R}}
\providecommand{\N}{\mathbb{N}}
\providecommand{\Z}{\mathbb{Z}}
\providecommand{\Q}{\mathbb{Q}}
\providecommand{\D}{\mathbb{D}}
\providecommand{\C}{\mathbb{C}}
\providecommand{\e}{{\mathrm{e}}}
\providecommand{\vect}[1]{\overrightarrow{#1}}
\providecommand{\vectt}[1]{\overrightarrow{\mathrm{#1}}}
\providecommand{\vvt}[1]{\overrightarrow{\mathrm{#1}}}
\providecommand{\np}[2][]{{#2~#1}}
\providecommand{\vv}[1]{{\overrightarrow{#1}}}
\renewcommand{\d}{{\mathrm{d}}}
\providecommand{\degres}{{^\circ}}
\providecommand{\degre}{{^\circ}}
\providecommand{\blue}{\color{blue}}
\providecommand{\red}{\color{red}}
\providecommand{\enskip}{\quad}
\providecommand{\llbracket}{\lbrack\!\lbrack}
\providecommand{\rrbracket}{\rbrack\!\rbrack}
\providecommand{\euro}{{\text{€}}}
\providecommand{\no}{{\text{n°}}}
\providecommand{\interval}[2]{{[#1;#2]}}
\providecommand{\intervFO}[2]{{[#1;#2[}}
\providecommand{\dotfill}{\ldots}

\providecommand{\ds}{\displaystyle}
\providecommand{\textding}[1]{#1}
\providecommand{\Oij}{\(\mathrm{O};\vec{\imath},\vec{\jmath}\)}
\providecommand{\Oijk}{\(\mathrm{O};\vec{\imath},\vec{\jmath},\vec{k}\)}
\providecommand{\Ouv}{\(\mathrm{O};\vec{u},\vec{v}\)}
\providecommand{\intervFF}[2]{\ensuremath{[#1;#2]}}
\renewcommand{\euro}{\ifmmode\text{\texteuro}\else\texteuro\fi}
\DeclareUnicodeCharacter{00B0}{\ifmmode^{\circ}\else\textdegree\fi}
\DeclareUnicodeCharacter{20AC}{\ifmmode\text{\texteuro}\else\texteuro\fi}
% Énoncés : Éducation nationale. Transcription et corrigés : APMEP, Denis Vergès et contributeurs. Découpage et classement : CoopMaths, Sébastien Lozano, Éric Elter et contributeurs. Adaptation de la présentation : Maths Explorer.
% CC BY-SA : https://coopmaths.fr/www/about/
% https://coopmaths.fr/alea/static/bac/2022/tex/bac_2022_05_sujet2_ameriquenord_2_cor.tex
\begin{document}
\noindent\textit{APMEP / CoopMaths. CC BY-SA. Présentation adaptée par Maths Explorer.}\par\medskip
\medskip

\textbf{Partie A}

\medskip

\begin{enumerate}
	\item La fonction $p$ est continue et dérivable sur $[-3~;~4]$.
	
	$\forall x \in [-3~;~4]$, $p'(x)= 3x^2-6x+5$
	
	Ce trinôme du second degré n'admet aucune racine ($\Delta=-24<0$) donc $\forall x \in [-3~;~4]$, $p'(x)>0$.
	Donc la fonction $p$ est strictement croissante sur $[-3~;~4]$.
	\item $p(-3)=-68$ et $p(4)=37$
	
	La fonction $p$ est continue et strictement croissante sur $[-3~;~4]$ à valeurs dans $[-68~;~37]$. Or $0\in [-68~;~37]$, donc d'après le corollaire du théorème des valeurs intermédiaires, l'équation $p(x)=0$ admet une unique solution, notée $\alpha$, dans l'intervalle $[-3~;~4]$.
	\item À la calculatrice, $\alpha \approx -0,2$.
	\item D'après les variations de la fonction $p$, et en utilisant le résultat précédent, on peut établir le tableau de signe de la fonction $p$ sur $[-3~;~4]$ :
	\begin{center}
	\psset{unit=1cm}
\begin{pspicture}(7,1.6)
\psframe(7,1.6)\psline(0,0.8)(7,0.8)\psline(2,0)(2,1.6)
\rput(1,1.2){$x$}\rput(1,0.4){$p(x)$}\rput(2.25,1.2){$-3$}\rput(4.5,1.2){$\alpha$}\rput(6.8,1.2){4}
\rput(3.25,0.4){$-$}\rput(4.5,0.4){0} \rput(5.75,0.4){$+$}
\end{pspicture}




	\end{center}

\end{enumerate}

\medskip

\textbf{Partie B}

\medskip


\begin{enumerate}
	\item 
	\begin{enumerate}
		\item La fonction $f$ est continue et dérivable sur $[-3~;~4]$ (car $\forall x \in [-3~;~4]$, $1+x^2\neq 0$).
		
		$\forall x \in [-3~;~4]$, $f'(x)=\dfrac{\e^x\times (1+x^2)-\e^x\times 2x}{\left(1+x^2\right)^2}=\dfrac{\e^x(x^2-2x+1)}{(1+^2)^2}=\dfrac{(x-1)^2~\e^x}{(1+^2)^2}$.
		\item $f'(x)=0\iff \dfrac{(x-1)^2~\e^x}{(1+^2)^2}=0 \iff (x-1)^2~\e^x \iff (x-1)^2 \iff x-1=0 \iff x=1$.
		
		Et $f(1)=\dfrac{\e}{2}$. Donc au point d'abscisse 1, $\mathcal{C}_f$ admet une tangente horizontale d'équation $y=\dfrac{\e}{2}$.
	\end{enumerate}
	\item
		\begin{enumerate}
			\item Avec la précision permise par le graphique, on peut voir que la fonction $f$ est :
				\begin{itemize}
					\item convexe sur $[-3~;~0]$ ;
					\item concave sur $[0~;~1]$ ;
					\item convexe sur $[1~;~4]$.
				\end{itemize}
		Donc $\mathcal{C}_f$ admet deux points d'inflexion, aux abscisses $x=0$ et $x=1$.
	
		Le toboggan semble donc assurer de bonnes sensations.
			\item $\forall x \in [-3~;~4]$ :\:$f''(x) = \dfrac{p(x)(x - 1)\text{e}^x}{\left(1 + x^2\right)^3}$
			
			Recherchons les points d'inflexion, c'est-à-dire les valeurs de $x\in [-3~;~4]$ pour lesquelles $f''(x)$ s'annule et change de signe.
			
			$\forall x \in [-3~;~4]$, $\left(1 + x^2\right)^3>0$ et $\e^x>0$ donc $f''(x)$ a le même signe que $p(x)(x - 1)$.
			
			On construit alors le tableau de signe suivant :
\begin{center}	
\psset{unit=1cm}
\begin{pspicture}(7,3.2)
\psframe(7,3.2)\psline(0,0.8)(7,0.8)\psline(0,1.6)(7,1.6)\psline(0,2.4)(7,2.4)\psline(1,0)(1,3.2)
\rput(0.5,2.8){$x$}\rput(1.25,2.8){$-3$}\rput(3,2.8){$\alpha$}\rput(5,2.8){$1$}\rput(6.8,2.8){$4$}
\rput(0.5,2){$p(x)$}\rput(2,2){$-$}\rput(3,2){$0$}\rput(4,2){$+$}\rput(6,2){$+$}
\rput(0.5,1.2){$x-1$}\rput(2,1.2){$-$}\rput(4,1.2){$-$}\rput(5,1.2){$0$}\rput(6,1.2){$+$}
\rput(0.5,0.4){$f''(x)$}\rput(2,0.4){$+$}\rput(3,0.4){$0$}\rput(4,0.4){$-$}\rput(5,0.4){$0$}\rput(6,0.4){$+$}
\end{pspicture}






			\end{center}
		La fonction $f''$ s'annule et change de signe en $x=\alpha$ et $x=1$. Donc $\mathcal{C}_f$ admet deux points d'inflexion. Le toboggan assure donc de bonnes sensations.
\end{enumerate} 
	
\end{enumerate}

\bigskip

\end{document}
